\section{Disconnected localization and quantitative path stability}
\label{sec:path-quantitative-localization}

Localization does not require the retained vertex set to be connected.
Removing that requirement avoids the component-counting step in the original stability proof.
The endpoint forest estimate then gives a quantitative additive stability constant.
All logarithms in this section have base two.

\begin{lemma}[Localization on a disconnected shore]
\label{lem:path-disconnected-localization}
Let $G$ be a connected multigraph, and let $R$ be a nonempty proper vertex set such that $G-R$ is connected.
Write
\[
 d=|\delta_G R|\ge1,\qquad c_R=\operatorname{comp}(G[R]).
\]
For each boundary edge, add a distinct new vertex and a spoke to its endpoint in $R$.
Join the new vertices into a cycle $W_R$, padding it to length three when $d<3$.
Let $G_R$ consist of $G[R]$, these spokes, and $W_R$.
Then $G_R$ is connected and
\[
 |V(W_R)|=\max\{d,3\},
 \qquad
 \beta(G_R)=\beta(G[R])+d-c_R+1.
\]
For uniform even subgraphs,
\[
 \Pr\bigl(\xi_G\cap E(G[R])\text{ is a forest}\bigr)
 =\Pr\bigl(\xi_{G_R}\cap E(G_R-V(W_R))\text{ is a forest}\bigr).
\]
If every vertex of $R$ has degree three in $G$, then
\[
 \beta(G_R)=\frac{|R|+d}{2}+1.
\]
\end{lemma}

\begin{proof}
Every component of $G[R]$ has a boundary edge, since $G$ is connected.
The new cycle therefore connects all these components.
Without padding, the construction adds $d$ vertices and $2d$ edges.
Its rank is $|E(G[R])|+d-|R|+1$, giving the stated formula.
Padding adds equally many edges and vertices and preserves this rank.

Contract all edges of the connected graph $G-R$ in $G$.
Separately, contract all edges of $W_R$ in $G_R$.
The resulting multigraphs agree, including their restrictions to $E(G[R])$.
Projection of a uniform cycle-space vector through a contraction is uniform on the image cycle space.
For surjectivity, complete each prescribed even boundary by a join in a spanning tree of the contracted component.
Thus the restrictions to $E(G[R])$ have the same law in both constructions.
This proves the probability identity.

Under the cubic assumption, $3|R|=2|E(G[R])|+d$.
Substitution into the rank formula gives $(|R|+d)/2+1$.
No connectedness assumption on $R$ was used.
\end{proof}

\begin{lemma}[A quantitative localized forest threshold]
\label{lem:path-polynomial-local-threshold}
There is an absolute constant $K_0$ with the following property.
For $0<\delta\le1/4$, one can choose an integer $R_\delta\ge2^{100}$ satisfying
\[
 R_\delta\le K_0\delta^{-2}\bigl(\log(2/\delta)\bigr)^{11}.
\]
Whenever a connected subcubic multigraph has rank $\rho\ge R_\delta$ and a connected marked set $P$ containing every vertex of degree below three, with
\[
 |P|\le\frac{\rho}{(\log\rho)^2},
\]
the probability that its uniform even subgraph is a forest outside $P$ is at most $1/2+\delta$.
\end{lemma}

\begin{proof}
Use the absolute constant $M$ and absolute rank threshold in Theorem~\ref{thm:parallel-links-log-five-half-forest}.
Increase that threshold to at least $2^{100}$.
At these ranks, a connected marked set satisfies its component bound.
It is enough to require
\[
 M\frac{(\log\rho)^{11/2}}{\sqrt\rho}\le\delta.
\]
The left side decreases for all sufficiently large $\rho$.
Choosing $\rho$ above a sufficiently large constant times
\[
 \delta^{-2}\bigl(\log(2/\delta)\bigr)^{11}
\]
ensures this inequality.
Indeed, the logarithm of this expression is $O(\log(2/\delta))$.
Choose the absolute constant large enough to cover both the endpoint threshold and the remaining compact range of $\delta$.
Rounding upward gives the asserted integer $R_\delta$.
\end{proof}

\begin{lemma}[Every sufficiently large good set is excluded]
\label{lem:path-direct-good-set}
Use the cubic expansion $\Gamma$, shortest terminal path $A$, graph $H$, and marked cycle $W$ from the general stability proof.
Let $\delta=\varepsilon/2$ and
\[
 \Theta=\max\{2^{36},2R_\delta\}.
\]
If the terminal-path density is at least $1/2+\varepsilon$, no set $R\subseteq V(\Gamma)\setminus V(A)$ can satisfy
\[
 |R|\ge\Theta,
 \qquad
 |\delta_\Gamma R|\le\frac{|R|}{4(\log|R|)^2}.
\]
\end{lemma}

\begin{proof}
Suppose such an $R$ exists, and put $m=|R|$.
Let $C_0$ be the component of $H-R$ containing $W$.
Set $R^+=V(H)\setminus V(C_0)$.
Then $R^+\supseteq R$, its complement is connected, and it avoids $V(W)$.
It need not be connected.
Every boundary edge of $R^+$ already belongs to $\delta_H R$, because distinct components of $H-R$ have no edge between them.
Hence, writing $d^+=|\delta_HR^+|$,
\[
 d^+\le|\delta_HR|=|\delta_\Gamma R|
 \le\frac{m}{4(\log m)^2}.
\]
All vertices of $R^+$ are cubic in $H$.

Apply Lemma~\ref{lem:path-disconnected-localization} to $R^+$.
The resulting graph has rank
\[
 \rho=\frac{|R^+|+d^+}{2}+1\ge m/2\ge R_\delta.
\]
It is subcubic, and every vertex outside its marked cycle is cubic.
The marked cycle is connected and has at most $\rho/(\log\rho)^2$ vertices.
For $d^+\ge3$, this follows from the increasing function $\Phi(t)=t/(\log t)^2$ on $t\ge8$ and
\[
 d^+\le\Phi(m)/4\le\Phi(m/2)\le\Phi(\rho).
\]
For $d^+<3$, padding adds only three marked vertices, which also meets that bound.

The terminal-path injection bounds $1/2+\varepsilon$ by the forest probability outside $W$ in $H$.
That event implies the forest event on $H[R^+]$.
The localization identity and Lemma~\ref{lem:path-polynomial-local-threshold} bound its probability by $1/2+\delta$.
This contradicts $\delta=\varepsilon/2$.
\end{proof}

\begin{theorem}[A single-exponential additive stability constant]
\label{thm:path-quantitative-additive}
There is an absolute constant $K_1$ such that the following holds.
Let $N$ be a connected graph of rank $b$ with distinct terminals $u,v$, and suppose
\[
 P_{uv}(N)\ge(1/2+\varepsilon)2^b,
 \qquad 0<\varepsilon\le1/2.
\]
Then every terminal path has length at least $b-C_\varepsilon$, where one may choose
\[
 C_\varepsilon\le
 2^{K_1\varepsilon^{-2}(\log(2/\varepsilon))^{11}}.
\]
The cubic expansion in the stability proof has distance layers of size at most
\[
 K_1\varepsilon^{-2}(\log(2/\varepsilon))^{11}+1.
\]
\end{theorem}

\begin{proof}
Lemma~\ref{lem:path-direct-good-set} replaces the original good-set dichotomy.
The unchanged growing-window lemma gives fewer than $\Theta$ unmarked vertices in each distance layer.
Lemma~\ref{lem:path-polynomial-local-threshold} bounds
\[
 \Theta\le K_2\varepsilon^{-2}(\log(2/\varepsilon))^{11}
\]
for an absolute $K_2$.
There is at most one marked path vertex in each layer.
This proves the layer-width statement after increasing the final constant.

Put $B=\lceil\Theta\rceil$ and $L=20B+2$.
The minimum-span window argument and greedy deletion of whole windows remain unchanged.
The already proved additive estimate therefore permits
\[
 C_\varepsilon=\frac{2L(L+1)2^{L-1}}{\varepsilon}.
\]
Taking logarithms, and using the displayed polynomial bound on $\Theta$, gives
\[
 \log C_\varepsilon
 \le K_1\varepsilon^{-2}(\log(2/\varepsilon))^{11}
\]
after increasing $K_1$.
The terminal reduction and rank-distance comparison then give the asserted path bound.
\end{proof}

\begin{corollary}[A quantitative edge-deletion path excess]
\label{cor:path-logarithmic-deletion-excess}
The function $q(s)$ from Proposition~\ref{prop:path-finite-cut-recurrence} satisfies
\[
 q(s)=O\!\left(\frac{(\log\log s)^{11/2}}{\sqrt{\log s}}\right).
\]
\end{corollary}

\begin{proof}
Put $x=\log s$, $z=\log x$, and choose a fixed $A\ge\max\{1,2\sqrt{K_1}\}$.
Set $\varepsilon=A z^{11/2}/\sqrt x$.
For sufficiently large $s$, this lies in $(0,1/2]$ and $\log(2/\varepsilon)\le z$.
Theorem~\ref{thm:path-quantitative-additive} then gives
\[
 C_\varepsilon
 \le2^{K_1\varepsilon^{-2}(\log(2/\varepsilon))^{11}}
 \le2^{x/4}=s^{1/4}.
\]
Consequently
\[
 B_{2,\varepsilon}=3(9\lceil C_\varepsilon\rceil+25)<s-1
\]
for large $s$.
The two-pole high-density small-cut theorem gives $q(s)\le\varepsilon$.
This proves the stated rate.
\end{proof}

\begin{corollary}[Growing cyclic cut orders]
\label{cor:path-loglog-growing-cut}
Let connected loopless cubic multigraphs $G$ have ranks $r\to\infty$ and minimum cyclic cut orders $k=k(r)\ge5$.
If
\[
 \frac{2^k(\log\log r)^{11/2}}{\sqrt{\log r}}\longrightarrow0,
\]
then
\[
 c(G)\le2^{r-1}-(1-o(1))2^{r-k-1}<2^{r-1}.
\]
In particular, this holds whenever
\[
 k\le(1/2-\eta)\log\log r
\]
for any fixed $\eta>0$.
\end{corollary}

\begin{proof}
The assumed condition gives $k=O(\log\log r)$, so $r-k\to\infty$.
Corollary~\ref{cor:path-logarithmic-deletion-excess} holds uniformly over the ranks from $r-k+6$ through $r$.
Its bound changes by a factor $1+o(1)$ across this interval.
Apply Proposition~\ref{prop:path-finite-cut-recurrence} and compare its error terms with $2^{-k-1}$.
Both are negligible under the displayed assumption.
For the final assertion, the ratio is at most
\[
 \frac{(\log\log r)^{11/2}}{(\log r)^\eta},
\]
which tends to zero.
\end{proof}

\begin{corollary}[The endpoint growing-order range]
\label{cor:path-endpoint-loglog-cut}
There is an absolute constant $C$ such that, for all sufficiently large $r$, every connected loopless cubic multigraph of rank $r$ with
\[
 5\le\lambda_c(G)=k
 \le\frac12\log\log r-\frac{11}{2}\log\log\log r-C
\]
satisfies
\[
 c(G)\le2^{r-1}-2^{r-k-2}.
\]
\end{corollary}

\begin{proof}
Choose an absolute $A_0$ such that the path-excess sum in Proposition~\ref{prop:path-finite-cut-recurrence} is at most
\[
 A_0\frac{(\log\log r)^{11/2}}{\sqrt{\log r}}
\]
throughout $k=O(\log\log r)$ at sufficiently large $r$.
This follows from Corollary~\ref{cor:path-logarithmic-deletion-excess} and the sum of its geometric weights, which is less than one.
Take $C$ so large that $A_0 2^{-C}\le1/8$.
The displayed cut-order bound then makes the path-excess sum at most $2^{-k-3}$.
For sufficiently large $r$, the base error satisfies $a_5(r-k+5)\le1/256$ uniformly in this range.
Its contribution is also at most $2^{-k-3}$.
Subtract these two errors from the deficit $2^{-k-1}$ in Proposition~\ref{prop:path-finite-cut-recurrence}.
The remaining deficit is at least $2^{-k-2}$, as asserted.
\end{proof}

The new localization step removes an unnecessary exponential loss in the stability constants.
It does not give the unrestricted half-exponent cycle error.
The remaining cyclic connectivity can still grow between the displayed double-logarithmic range and the general logarithmic upper bound.
